% preamble.tex: packages, macros and boxes for the study guide.
% Chapters use only the macros and environments defined here. Do not add packages in a chapter.

\usepackage[T1]{fontenc}
\usepackage[utf8]{inputenc}
\usepackage[a4paper,margin=1in,headheight=14pt]{geometry}
\usepackage{amsmath,amsthm,mathtools}
\usepackage{stix2}
\usepackage{microtype}
\usepackage{graphicx}
\usepackage{booktabs}
\usepackage{array}
\usepackage{enumitem}
\usepackage{needspace}
\usepackage[font=small,labelfont=bf,labelsep=period]{caption}
\usepackage{xcolor}
\usepackage{tikz}
\usetikzlibrary{arrows.meta,positioning,calc,decorations.pathreplacing}
\usepackage[most]{tcolorbox}
\usepackage{titlesec}
\usepackage{fancyhdr}
\usepackage{hyperref}

\graphicspath{{source/figures/}}

% ---------------------------------------------------------------- colours
% Grey ink and one dark accent. No bright colours.
\definecolor{ink}{HTML}{1F2328}
\definecolor{accent}{HTML}{2E4A62}
\definecolor{rulegrey}{HTML}{8C959F}
\definecolor{fillgrey}{HTML}{F4F5F6}
\definecolor{notegrey}{HTML}{57606A}
\color{ink}

\hypersetup{
  pdftitle={Competition Creates Competition: A Study Guide},
  pdfauthor={Austin Li},
  colorlinks=true,
  linkcolor=accent,
  citecolor=accent,
  urlcolor=accent}

% ---------------------------------------------------------------- layout
\setlength{\parindent}{0pt}
\setlength{\parskip}{0.6em}
\linespread{1.08}
\allowdisplaybreaks
\setlist{itemsep=0.25em,topsep=0.4em}

\titleformat{\chapter}[display]
  {\normalfont\bfseries\color{ink}}
  {\large\color{notegrey}\MakeUppercase{\chaptertitlename}~\thechapter}
  {0.4em}
  {\LARGE}
  [\vspace{0.3em}{\color{rulegrey}\titlerule[0.6pt]}]
\titlespacing*{\chapter}{0pt}{0pt}{2.2em}
\titleformat{\section}{\normalfont\large\bfseries\color{ink}}{\thesection}{0.8em}{}
\titleformat{\subsection}{\normalfont\normalsize\bfseries\color{ink}}{\thesubsection}{0.8em}{}

\pagestyle{fancy}
\fancyhf{}
\fancyhead[L]{\small\color{notegrey}\nouppercase{\leftmark}}
\fancyhead[R]{\small\color{notegrey}\thepage}
\renewcommand{\headrulewidth}{0.4pt}
\renewcommand{\headrule}{{\color{rulegrey}\hrule height\headrulewidth width\headwidth}}
\fancypagestyle{plain}{\fancyhf{}\fancyfoot[C]{\small\color{notegrey}\thepage}\renewcommand{\headrulewidth}{0pt}}

% ---------------------------------------------------------------- manuscript macros
% These match the manuscript, so formulas read the same in both documents.
\newcommand{\E}{\mathbb{E}}
\newcommand{\1}{\mathbf{1}}
\newcommand{\DT}{\Delta_T}
\newcommand{\eu}{\mathrm{e}}          % Euler's number (entry is e(x))
\newcommand{\Ent}{\mathsf{E}}          % entry, the preparation probability
\newcommand{\Own}{\mathsf{O}}          % ownership
\newcommand{\fr}{\mathfrak{r}}         % root of Delta_T(r) = d above ell
\newcommand{\lb}{\bar\lambda}          % research cap
\DeclareMathOperator{\Cov}{Cov}
\newcommand{\Prob}{\Pr}

% ---------------------------------------------------------------- derivation notes
% \why{rule} puts a short grey note at the right of a derivation line.
% Use it inside align or align*:   &= a+b \why{linearity of expectation}
% The note is set as a tag, so amsmath moves it to its own line when the equation is too wide.
\newcommand{\why}[1]{\tag*{\small\color{notegrey}[#1]}}
% \whyinline{rule} gives the same note inside running text or a single display.
\newcommand{\whyinline}[1]{\quad\text{\small\color{notegrey}[#1]}}

% ---------------------------------------------------------------- boxes
\tcbset{
  guidebox/.style={
    enhanced, breakable,
    colback=fillgrey, colframe=rulegrey,
    boxrule=0pt, leftrule=1.6pt, arc=0pt, outer arc=0pt,
    left=8pt, right=8pt, top=6pt, bottom=6pt,
    before skip=0.9em, after skip=0.9em,
    fonttitle=\bfseries\small, coltitle=accent,
    attach title to upper, after title={\quad}}
}

% In this chapter: the question and the plan. One per chapter, at the start.
\newtcolorbox{chapterplan}{guidebox, colback=white, leftrule=0pt, toprule=0.6pt, bottomrule=0.6pt,
  colframe=rulegrey, title={In this chapter}, attach title to upper=false, detach title, before upper={\tcbtitle\par\smallskip}}

% Intuition in words, before or after the algebra. Optional title.
\newtcolorbox{intuition}[1][]{guidebox, title={In words\ifx\relax#1\relax\else: #1\fi}}

% One-sentence key idea.
\newtcolorbox{keyidea}{guidebox, leftrule=3pt, title={Key idea}}

% A result of the paper, stated with its exact hypotheses. Argument: the name.
\newtcolorbox{paperresult}[1]{guidebox, colback=white, boxrule=0.6pt, leftrule=0.6pt,
  title={#1}, attach title to upper=false, detach title, before upper={\tcbtitle\par\smallskip\itshape}}

% A hypothesis or a distinction that is easy to lose.
\newtcolorbox{watchout}{guidebox, leftrule=1.6pt, colframe=ink, title={Watch out}}

% A question for the reader. The answer goes in the chapter's Answers section.
\newcounter{checkq}[chapter]
\renewcommand{\thecheckq}{\thechapter.\arabic{checkq}}
% The counter steps outside the box, so a \label inside the box takes the question number.
\newtcolorbox{checkyourselfbox}{guidebox, colback=white, leftrule=1.6pt, colframe=accent,
  title={Check yourself \thecheckq}}
\newenvironment{checkyourself}{\refstepcounter{checkq}\begin{checkyourselfbox}}{\end{checkyourselfbox}}

% An analogy with its limit. Optional title.
\newtcolorbox{analogy}[1][]{guidebox, title={An analogy\ifx\relax#1\relax\else: #1\fi}}

% A derivation block: no fill, a thin left rule, can break across pages.
\newtcolorbox{derivation}[1][]{enhanced, breakable, colback=white, colframe=rulegrey,
  boxrule=0pt, leftrule=0.8pt, arc=0pt, outer arc=0pt, left=8pt, right=0pt, top=2pt, bottom=2pt,
  before skip=0.8em, after skip=0.8em,
  title={Derivation\ifx\relax#1\relax\else: #1\fi}, fonttitle=\bfseries\small, coltitle=accent,
  attach title to upper, after title={\par\smallskip}}

% Chapter summary at the end of the chapter.
\newtcolorbox{chaptersummary}{guidebox, colback=white, leftrule=0pt, toprule=0.6pt, bottomrule=0.6pt,
  title={What to remember}, attach title to upper=false, detach title, before upper={\tcbtitle\par\smallskip}}

% Answers section heading for the check-yourself questions.
\newcommand{\answersection}{\section*{Answers to the check-yourself questions}\addcontentsline{toc}{section}{Answers to the check-yourself questions}}
% \answer{3.2}{text} prints one worked answer.
\newcommand{\answer}[2]{\par\needspace{3\baselineskip}\textbf{\color{accent}#1.}\quad #2\par}
